\begin{lemma}
\label{lemma-exact-on-fibres-open}
Let $R \to S$ be a ring map. Consider a finite homological complex of
finite free $S$-modules:
$$
F_{\bullet} :
0
\to
S^{n_e}
\xrightarrow{\varphi_e}
S^{n_{e-1}}
\xrightarrow{\varphi_{e-1}}
\ldots
\xrightarrow{\varphi_{i + 1}}
S^{n_i}
\xrightarrow{\varphi_i}
S^{n_{i-1}}
\xrightarrow{\varphi_{i-1}}
\ldots
\xrightarrow{\varphi_1}
S^{n_0}
$$
For every prime $\mathfrak q$ of $S$ consider the
complex $\overline{F}_{\bullet, \mathfrak q} =
F_{\bullet, \mathfrak q} \otimes_R \kappa(\mathfrak p)$
where $\mathfrak p$ is inverse image of $\mathfrak q$ in $R$.
Assume $R$ is Noetherian and there exists an integer $d$ such
that $R \to S$ is finite type, flat
with fibres $S \otimes_R \kappa(\mathfrak p)$
Cohen-Macaulay of dimension $d$.
The set
$$
\{\mathfrak q \in \Spec(S) \mid
\overline{F}_{\bullet, \mathfrak q}\text{ is exact}\}
$$
is open in $\Spec(S)$.
\end{lemma}

\begin{proof}
Let $\mathfrak q$ be an element of the set defined in the lemma.
We are going to use Proposition \ref{proposition-what-exact}
to show there exists a $g \in S$, $g \not \in \mathfrak q$
such that $D(g)$ is contained in the set defined in the lemma.
In other words, we are going to show that after replacing $S$
by $S_g$, the set of the lemma is all of $\Spec(S)$.
Thus during the proof we will, finitely often, replace
$S$ by such a localization.
Recall that Proposition \ref{proposition-what-exact}
characterizes exactness of complexes
in terms of ranks of the maps $\varphi_i$ and the ideals
$I(\varphi_i)$, in case the ring is local. We first address
the rank condition. Set
$r_i = n_i - n_{i + 1} + \ldots + (-1)^{e - i} n_e$.
Note that $r_i + r_{i + 1} = n_i$ and note that
$r_i$ is the expected rank of $\varphi_i$ (in the
exact case).

\medskip\noindent
By Lemma \ref{lemma-complex-exact-mod} we see that if
$\overline{F}_{\bullet, \mathfrak q}$ is exact, then
the localization $F_{\bullet, \mathfrak q}$ is exact.
In particular the complex $F_\bullet$ becomes
exact after localizing by an element
$g \in S$, $g \not \in \mathfrak q$. In this case
Proposition \ref{proposition-what-exact} applied
to all localizations of $S$ at prime ideals
implies that all $(r_i + 1) \times (r_i + 1)$-minors
of $\varphi_i$ are zero. Thus we see that the rank
of $\varphi_i$ is at most $r_i$.

\medskip\noindent
Let $I_i \subset S$ denote the ideal generated
by the $r_i \times r_i$-minors of the matrix
of $\varphi_i$. By Proposition \ref{proposition-what-exact}
the complex $\overline{F}_{\bullet, \mathfrak q}$ is exact
if and only if for every $1 \leq i \leq e$ we have
either $(I_i)_{\mathfrak q} = S_{\mathfrak q}$ or
$(I_i)_{\mathfrak q}$ contains a $S_{\mathfrak q}/\mathfrak p
S_{\mathfrak q}$-regular sequence of length $i$.
Namely, by our choice of $r_i$ above and by the
bound on the ranks of the $\varphi_i$ this is the
only way the conditions of Proposition \ref{proposition-what-exact}
can be satisfied.

\medskip\noindent
If $(I_i)_{\mathfrak q} = S_{\mathfrak q}$, then after localizing $S$ at
some element $g \not\in \mathfrak q$ we may assume that
$I_i = S$. Clearly, this is an open condition.

\medskip\noindent
If $(I_i)_{\mathfrak q} \not = S_{\mathfrak q}$, then we have
a sequence $f_1, \ldots, f_i \in (I_i)_{\mathfrak q}$ which
form a regular sequence in $S_{\mathfrak q}/\mathfrak pS_{\mathfrak q}$.
Note that for any prime $\mathfrak q' \subset S$ such that
$(f_1, \ldots, f_i) \not \subset \mathfrak q'$ we have
$(I_i)_{\mathfrak q'} = S_{\mathfrak q'}$.
Thus the result follows from Lemma \ref{lemma-open-regular-sequence}.
\end{proof}